\section{Resumen ejecutivo}

Este proyecto propone un sistema distribuido para gestionar los documentos y
las actividades de las pre-sustentaciones de trabajos de titulación de la
carrera de Software de la Universidad Técnica Estatal de Quevedo. Actualmente,
parte del proceso se coordina de manera manual y por correo electrónico. Las
actas se solicitan y entregan por ese medio, las firmas en papel retrasan el
cierre y cambios importantes, como la sustitución de un tutor, pueden quedar
sin registro. La entrevista realizada a la Ing. Nancy Rodríguez permitió
precisar que normalmente participan entre 10 y 12 estudiantes por periodo y
que la calificación, el acta y el porcentaje de similitud son datos críticos.

La propuesta cubre el registro y seguimiento de solicitudes, la entrega de
informes en PDF y reportes de URKUND, la asignación de tutores y tribunales, el
registro de evaluaciones, la consolidación de la nota y la generación del acta.
Se definieron doce historias de usuario y requisitos no funcionales orientados
a rendimiento, disponibilidad, seguridad, integridad y facilidad de uso. La
nota final combina el 40\% del tribunal con el 60\% del docente de Desarrollo
de Titulación II; por ello, las calificaciones y las actas requieren
consistencia, mientras que avisos, listados y reportes admiten una actualización
no inmediata.

La arquitectura se documentó mediante los niveles de contexto y contenedores
del modelo C4. Los usuarios acceden desde una aplicación web Angular a través
de un API Gateway. La solución separa cinco servicios Spring Boot: Usuarios,
Cronograma, Solicitudes y documentos, Evaluación y actas, y Notificaciones.
Cada servicio administra su propia base PostgreSQL; los documentos se almacenan
en MinIO y RabbitMQ distribuye los eventos que no necesitan respuesta
inmediata. Esta separación busca que una falla en notificaciones o cronograma
no impida registrar evaluaciones durante la semana de predefensas. Los ADR
justifican el estilo de microservicios y la combinación de comunicación
síncrona y asíncrona.

El análisis contempla riesgos de interrupción en la carga de documentos,
lecturas de notas desactualizadas, caída de componentes compartidos, cambios
defectuosos en reglas y saturación de la API. Las mitigaciones incluyen
validación de archivos, lectura de la fuente principal para generar actas,
despliegues graduales, versionado de reglas, acceso administrativo independiente
y límites de solicitudes. Como resultado de esta primera entrega quedan
establecidos el problema, los requisitos, la arquitectura base, las decisiones
técnicas, los riesgos y un esqueleto funcional de los servicios principales.
